\section{Introduction}
\label{sec:introduction}

The completion problem now has an exact analytic optimum, but a physical
rule must remember where its coefficients came from.  For one actual
row pair, TPC--74 \citep{WangTPC74} gives
\[
 L(u,v)=\sum_{q\mid|\Delta|}L_q(u,v),
 \qquad
 \Delta=h_0(n-m)\ne0,
\]
on the primitive cofactor plane.  The \(q\)-slice is not an auxiliary
Fourier label: it is the greatest common content of the two literal
residual factors, and its divisor restriction follows from the exact
fixed-shift Diophantine identity.  Aggregation over \(q\) must occur
before an absolute-value estimate, because distinct slices can cancel at
the same primitive coordinate.

TPC--75 \citep{WangTPC75} proves that, for a finite rectangle and full
divisor--variation analysis operator \(A_D\),
\[
 \gamma_D(L)
 =
 \min_{RF=L}\|A_DF\|_1
\]
has an exact strong dual.  That theorem deliberately allows arbitrary
nonprimitive values.  Such freedom provides the sharpest separated
analytic benchmark, but it may mix literal contents in a way no uniform
physical rule can reproduce.

The present paper inserts an intermediate class:
\[
 \text{zero completion}
 \subset
 \text{content-filter completions}
 \subset
 \text{all analytic completions}.
\]
The class is generated only from the actual \(q\)-resolved kernels.  It
therefore preserves fixed-\(h_0\) provenance without declaring its
nonprimitive values to be new atoms.

\subsection{The divisor-poset lift}

Every lattice point has a unique representation
\[
 (x,y)=(du,dv),
 \qquad d=(x,y),\quad(u,v)=1.
\]
We fill that point with the cumulative content above \(d\):
\[
 F^{\rm Z}(du,dv)
 =
 \sum_{d\mid q\mid|\Delta|}L_q(u,v).
\]
This is the upper zeta transform on the finite divisor poset of
\(|\Delta|\).  Poset M\"obius inversion \citep{Rota1964} recovers every
literal slice.  Hence the lift is information-preserving on the full
divisor profile even though its primitive trace sees only the aggregate.

The same cumulative structure identifies every diagonal slice exactly.
Only \(a\mid|\Delta|\) can contribute to the \(a\)-th divisor lift.
This is a genuine carrier reduction relative to an arbitrary ambient
completion.  It is not a tail estimate: a function supported on few
diagonal indices can still have large coefficients.

\subsection{A provenance-compatible optimization}

The zeta coefficient \(1\) is only one rule.  Allow weights
\(w_{d,q}\) whenever \(d\mid q\), while fixing \(w_{1,q}=1\) so the
primitive datum is unchanged.  Both the zero completion and zeta lift
belong to this affine family.  For fixed literal slices, the completion is
linear in the free weights, so its exact TPC--75 cost becomes a smaller
finite convex program.  We derive its reduced strong dual.  The dual
stationarity condition says that the analysis charge annihilates every
permitted content-flow direction.

\subsection{What is and is not gained}

The new results are exact:
\begin{itemize}
 \item the lift, inverse, diagonal support, and Hermitian reversal;
 \item a sharp support census in terms of the actual \(q\)-slice supports;
 \item the full provenance-filter hierarchy;
 \item primal and dual attainment in that hierarchy; and
 \item exact operator upper and obstruction matrices after row-pair
       aggregation.
\end{itemize}
They do not order the zero and zeta costs universally.  Lower
\(d=1\) variation can be exchanged for higher diagonal mass.  They also
do not turn the divisor bound
\(\#\{d:d\mid\Delta\}=\tau(|\Delta|)\) into a normalized saving.

\subsection{Organization}

\Cref{sec:zeta-lift} proves the zeta lift and inversion.
\Cref{sec:diagonal-filtration} derives the exact diagonal hierarchy and
support census.  \Cref{sec:provenance-gauge} introduces all content
filters and their reduced dual.  \Cref{sec:Hermitian-aggregation}
handles reversal and full matrices.  \Cref{sec:countermodels} records
the non-substitution laws.  \Cref{sec:physical-gate} states the actual
fixed-shift gate, and \cref{sec:claim-boundary} gives the audit ledger.
